\documentclass[letterpaper]{article}
\usepackage[submission]{aaai2027}
\usepackage[hyphens]{url}
\usepackage{graphicx}
\urlstyle{rm}
\def\UrlFont{\rm}
\usepackage{natbib}
\usepackage{caption}
\usepackage{booktabs}
\usepackage{amsmath}
\frenchspacing

\pdfinfo{
/TemplateVersion (2027.1)
/Title (LeakBench-Tab: Separating Contamination Validity, Statistical Detectability, and Model Exploitability in Tabular Learning)
/Author (Anonymous Submission)
/Subject (AAAI-27 Main Technical Track Submission)
}

\setcounter{secnumdepth}{2}
% Safe defaults used until generated/result_macros.tex is produced from the
% complete confirmatory paper_claims.json.  The front-page draft banner is
% controlled by \ifLBResultsReady in main.tex and supplement.tex.
\newif\ifLBResultsReady
\newif\ifLBSimpleStructuredSupported
\LBResultsReadyfalse
\LBSimpleStructuredSupportedfalse

\newcommand{\LBPending}{\ensuremath{\mathrm{PENDING}}}
\newcommand{\LBSourceSHA}{\LBPending}
\newcommand{\LBCDXScatterPath}{figures/generated/cdx_scatter.pdf}
\newcommand{\LBMechanismModelHeatmapPath}{figures/generated/mechanism_model_heatmap.pdf}
\newcommand{\LBStrengthDiagnosticPath}{figures/generated/strength_diagnostic_robustness.pdf}
\newcommand{\LBControlledTaskCount}{20}
\newcommand{\LBMechanismCount}{11}
\newcommand{\LBStrengthCount}{5}
\newcommand{\LBCoreModelCount}{5}
\newcommand{\LBSeedCount}{5}
\newcommand{\LBExpectedCells}{27,500}
\newcommand{\LBSuccessfulCells}{\LBPending}
\newcommand{\LBCompletionRate}{\LBPending}
\newcommand{\LBDiagnosticMethodCount}{4}
\newcommand{\LBExpectedDiagnosticCells}{22,000}
\newcommand{\LBSuccessfulDiagnosticCells}{\LBPending}
\newcommand{\LBDiagnosticCompletionRate}{\LBPending}
\newcommand{\LBDiagnosticConditionalStatus}{pending}
\newcommand{\LBSimpleStructuredStatus}{pending}
\newcommand{\LBSimpleStructuredDifference}{\LBPending}
\newcommand{\LBSimpleStructuredCILow}{\LBPending}
\newcommand{\LBSimpleStructuredCIHigh}{\LBPending}
\newcommand{\LBSimpleStructuredHolmP}{\LBPending}
\newcommand{\LBModelPositiveDirectionCount}{\LBPending}
\newcommand{\LBModelCIExcludesZeroCount}{\LBPending}

\newcommand{\LBMThreeStatus}{pending}
\newcommand{\LBMThreeDetectability}{\LBPending}
\newcommand{\LBMThreeDetectabilityCILow}{\LBPending}
\newcommand{\LBMThreeDetectabilityCIHigh}{\LBPending}
\newcommand{\LBMThreeHarm}{\LBPending}
\newcommand{\LBMThreeHarmCILow}{\LBPending}
\newcommand{\LBMThreeHarmCIHigh}{\LBPending}

\newcommand{\LBMEightStatus}{pending}
\newcommand{\LBMEightDetectability}{\LBPending}
\newcommand{\LBMEightDetectabilityCILow}{\LBPending}
\newcommand{\LBMEightDetectabilityCIHigh}{\LBPending}
\newcommand{\LBMEightHarm}{\LBPending}
\newcommand{\LBMEightHarmCILow}{\LBPending}
\newcommand{\LBMEightHarmCIHigh}{\LBPending}
\newcommand{\LBMEightClusterCILow}{\LBPending}
\newcommand{\LBMEightClusterCIHigh}{\LBPending}

\newcommand{\LBMNineStatus}{pending}
\newcommand{\LBMNineDetectability}{\LBPending}
\newcommand{\LBMNineDetectabilityCILow}{\LBPending}
\newcommand{\LBMNineDetectabilityCIHigh}{\LBPending}
\newcommand{\LBMNineHarm}{\LBPending}
\newcommand{\LBMNineHarmCILow}{\LBPending}
\newcommand{\LBMNineHarmCIHigh}{\LBPending}
\newcommand{\LBMNineReweightingLow}{\LBPending}
\newcommand{\LBMNineReweightingHigh}{\LBPending}

\newcommand{\LBRelationStatus}{pending}
\newcommand{\LBGlobalSpearman}{\LBPending}
\newcommand{\LBGlobalSpearmanCILow}{\LBPending}
\newcommand{\LBGlobalSpearmanCIHigh}{\LBPending}
\newcommand{\LBCategoryRSquared}{\LBPending}
\newcommand{\LBCategoryPlusDRSquared}{\LBPending}
\newcommand{\LBIncrementalRSquared}{\LBPending}
\newcommand{\LBIncrementalPermutationP}{\LBPending}
\newcommand{\LBCategoryLomoRSquared}{\LBPending}
\newcommand{\LBCategoryPlusDLomoRSquared}{\LBPending}
\newcommand{\LBIncrementalLomoRSquared}{\LBPending}

\newcommand{\LBNaturalStatus}{pending}
\newcommand{\LBNaturalTaskCount}{5}
\newcommand{\LBNaturalModelCount}{4}
\newcommand{\LBNaturalAllPositive}{\LBPending}
\newcommand{\LBNaturalMeanHarm}{\LBPending}
\newcommand{\LBNaturalCILow}{\LBPending}
\newcommand{\LBNaturalCIHigh}{\LBPending}
\newcommand{\LBNaturalSignFlipP}{\LBPending}

\newcommand{\LBMThreeMIDetectability}{\LBPending}
\newcommand{\LBMThreeMICILow}{\LBPending}
\newcommand{\LBMThreeMICIHigh}{\LBPending}
\newcommand{\LBMThreeCorrelationDetectability}{\LBPending}
\newcommand{\LBMThreeCorrelationCILow}{\LBPending}
\newcommand{\LBMThreeCorrelationCIHigh}{\LBPending}
\newcommand{\LBMThreeLRCoefficientDetectability}{\LBPending}
\newcommand{\LBMThreeLRCoefficientCILow}{\LBPending}
\newcommand{\LBMThreeLRCoefficientCIHigh}{\LBPending}
\newcommand{\LBMThreeRFPermutationDetectability}{\LBPending}
\newcommand{\LBMThreeRFPermutationCILow}{\LBPending}
\newcommand{\LBMThreeRFPermutationCIHigh}{\LBPending}
\newcommand{\LBMThreeDiagnosticMin}{\LBPending}
\newcommand{\LBMThreeDiagnosticMax}{\LBPending}

\newcommand{\LBMFourMIDetectability}{\LBPending}
\newcommand{\LBMFourMICILow}{\LBPending}
\newcommand{\LBMFourMICIHigh}{\LBPending}
\newcommand{\LBMFourCorrelationDetectability}{\LBPending}
\newcommand{\LBMFourCorrelationCILow}{\LBPending}
\newcommand{\LBMFourCorrelationCIHigh}{\LBPending}
\newcommand{\LBMFourLRCoefficientDetectability}{\LBPending}
\newcommand{\LBMFourLRCoefficientCILow}{\LBPending}
\newcommand{\LBMFourLRCoefficientCIHigh}{\LBPending}
\newcommand{\LBMFourRFPermutationDetectability}{\LBPending}
\newcommand{\LBMFourRFPermutationCILow}{\LBPending}
\newcommand{\LBMFourRFPermutationCIHigh}{\LBPending}
\newcommand{\LBMFourDiagnosticMin}{\LBPending}
\newcommand{\LBMFourDiagnosticMax}{\LBPending}

\newcommand{\LBMFiveMIDetectability}{\LBPending}
\newcommand{\LBMFiveMICILow}{\LBPending}
\newcommand{\LBMFiveMICIHigh}{\LBPending}
\newcommand{\LBMFiveCorrelationDetectability}{\LBPending}
\newcommand{\LBMFiveCorrelationCILow}{\LBPending}
\newcommand{\LBMFiveCorrelationCIHigh}{\LBPending}
\newcommand{\LBMFiveLRCoefficientDetectability}{\LBPending}
\newcommand{\LBMFiveLRCoefficientCILow}{\LBPending}
\newcommand{\LBMFiveLRCoefficientCIHigh}{\LBPending}
\newcommand{\LBMFiveRFPermutationDetectability}{\LBPending}
\newcommand{\LBMFiveRFPermutationCILow}{\LBPending}
\newcommand{\LBMFiveRFPermutationCIHigh}{\LBPending}
\newcommand{\LBMFiveDiagnosticMin}{\LBPending}
\newcommand{\LBMFiveDiagnosticMax}{\LBPending}

\IfFileExists{generated/result_macros.tex}{% AUTO-GENERATED by source_data/generate_result_macros.py.
% Do not edit by hand.
% paper_claims.json sha256: deff6fe15791439d16c64673061091eed62b5fc4f0f6bc60780fe7b5dfb64c04
\LBResultsReadytrue
\LBSimpleStructuredSupportedtrue
\renewcommand{\LBSourceSHA}{deff6fe15791439d16c64673061091eed62b5fc4f0f6bc60780fe7b5dfb64c04}
\renewcommand{\LBControlledTaskCount}{20}
\renewcommand{\LBMechanismCount}{11}
\renewcommand{\LBStrengthCount}{5}
\renewcommand{\LBCoreModelCount}{5}
\renewcommand{\LBSeedCount}{5}
\renewcommand{\LBExpectedCells}{27,500}
\renewcommand{\LBSuccessfulCells}{27,500}
\renewcommand{\LBCompletionRate}{100.0}
\renewcommand{\LBDiagnosticMethodCount}{4}
\renewcommand{\LBExpectedDiagnosticCells}{22,000}
\renewcommand{\LBSuccessfulDiagnosticCells}{22,000}
\renewcommand{\LBDiagnosticCompletionRate}{100.0}
\renewcommand{\LBDiagnosticConditionalStatus}{descriptive only}
\renewcommand{\LBSimpleStructuredStatus}{supported}
\renewcommand{\LBSimpleStructuredDifference}{0.164}
\renewcommand{\LBSimpleStructuredCILow}{0.150}
\renewcommand{\LBSimpleStructuredCIHigh}{0.178}
\renewcommand{\LBSimpleStructuredHolmP}{<0.001}
\renewcommand{\LBModelPositiveDirectionCount}{5}
\renewcommand{\LBModelCIExcludesZeroCount}{0}
\renewcommand{\LBMThreeStatus}{descriptive only}
\renewcommand{\LBMThreeDetectability}{0.079}
\renewcommand{\LBMThreeDetectabilityCILow}{0.045}
\renewcommand{\LBMThreeDetectabilityCIHigh}{0.128}
\renewcommand{\LBMThreeHarm}{0.217}
\renewcommand{\LBMThreeHarmCILow}{0.201}
\renewcommand{\LBMThreeHarmCIHigh}{0.234}
\renewcommand{\LBMEightStatus}{descriptive only}
\renewcommand{\LBMEightDetectability}{0.356}
\renewcommand{\LBMEightDetectabilityCILow}{0.263}
\renewcommand{\LBMEightDetectabilityCIHigh}{0.457}
\renewcommand{\LBMEightHarm}{0.003}
\renewcommand{\LBMEightHarmCILow}{-0.002}
\renewcommand{\LBMEightHarmCIHigh}{0.008}
\renewcommand{\LBMNineStatus}{descriptive only}
\renewcommand{\LBMNineDetectability}{0.541}
\renewcommand{\LBMNineDetectabilityCILow}{0.500}
\renewcommand{\LBMNineDetectabilityCIHigh}{0.583}
\renewcommand{\LBMNineHarm}{0.235}
\renewcommand{\LBMNineHarmCILow}{0.216}
\renewcommand{\LBMNineHarmCIHigh}{0.254}
\renewcommand{\LBMEightClusterCILow}{-0.006}
\renewcommand{\LBMEightClusterCIHigh}{0.010}
\renewcommand{\LBMNineReweightingLow}{0.199}
\renewcommand{\LBMNineReweightingHigh}{0.235}
\renewcommand{\LBRelationStatus}{descriptive only}
\renewcommand{\LBGlobalSpearman}{0.555}
\renewcommand{\LBGlobalSpearmanCILow}{0.482}
\renewcommand{\LBGlobalSpearmanCIHigh}{0.709}
\renewcommand{\LBCategoryRSquared}{0.514}
\renewcommand{\LBCategoryPlusDRSquared}{0.586}
\renewcommand{\LBIncrementalRSquared}{0.072}
\renewcommand{\LBIncrementalPermutationP}{0.330}
\renewcommand{\LBCategoryLomoRSquared}{0.086}
\renewcommand{\LBCategoryPlusDLomoRSquared}{-0.993}
\renewcommand{\LBIncrementalLomoRSquared}{-1.079}
\renewcommand{\LBNaturalStatus}{case-study only}
\renewcommand{\LBNaturalTaskCount}{5}
\renewcommand{\LBNaturalModelCount}{4}
\renewcommand{\LBNaturalAllPositive}{yes}
\renewcommand{\LBNaturalMeanHarm}{0.273}
\renewcommand{\LBNaturalCILow}{0.120}
\renewcommand{\LBNaturalCIHigh}{0.424}
\renewcommand{\LBNaturalSignFlipP}{0.062}
\renewcommand{\LBMThreeMIDetectability}{0.079}
\renewcommand{\LBMThreeMICILow}{0.069}
\renewcommand{\LBMThreeMICIHigh}{0.089}
\renewcommand{\LBMThreeCorrelationDetectability}{0.035}
\renewcommand{\LBMThreeCorrelationCILow}{0.025}
\renewcommand{\LBMThreeCorrelationCIHigh}{0.045}
\renewcommand{\LBMThreeLRCoefficientDetectability}{0.045}
\renewcommand{\LBMThreeLRCoefficientCILow}{0.035}
\renewcommand{\LBMThreeLRCoefficientCIHigh}{0.055}
\renewcommand{\LBMThreeRFPermutationDetectability}{0.602}
\renewcommand{\LBMThreeRFPermutationCILow}{0.592}
\renewcommand{\LBMThreeRFPermutationCIHigh}{0.612}
\renewcommand{\LBMThreeDiagnosticMin}{0.035}
\renewcommand{\LBMThreeDiagnosticMax}{0.602}
\renewcommand{\LBMFourMIDetectability}{0.081}
\renewcommand{\LBMFourMICILow}{0.071}
\renewcommand{\LBMFourMICIHigh}{0.091}
\renewcommand{\LBMFourCorrelationDetectability}{0.039}
\renewcommand{\LBMFourCorrelationCILow}{0.029}
\renewcommand{\LBMFourCorrelationCIHigh}{0.049}
\renewcommand{\LBMFourLRCoefficientDetectability}{0.037}
\renewcommand{\LBMFourLRCoefficientCILow}{0.027}
\renewcommand{\LBMFourLRCoefficientCIHigh}{0.047}
\renewcommand{\LBMFourRFPermutationDetectability}{0.045}
\renewcommand{\LBMFourRFPermutationCILow}{0.035}
\renewcommand{\LBMFourRFPermutationCIHigh}{0.055}
\renewcommand{\LBMFourDiagnosticMin}{0.037}
\renewcommand{\LBMFourDiagnosticMax}{0.081}
\renewcommand{\LBMFiveMIDetectability}{0.084}
\renewcommand{\LBMFiveMICILow}{0.074}
\renewcommand{\LBMFiveMICIHigh}{0.094}
\renewcommand{\LBMFiveCorrelationDetectability}{0.037}
\renewcommand{\LBMFiveCorrelationCILow}{0.027}
\renewcommand{\LBMFiveCorrelationCIHigh}{0.047}
\renewcommand{\LBMFiveLRCoefficientDetectability}{0.036}
\renewcommand{\LBMFiveLRCoefficientCILow}{0.026}
\renewcommand{\LBMFiveLRCoefficientCIHigh}{0.046}
\renewcommand{\LBMFiveRFPermutationDetectability}{0.044}
\renewcommand{\LBMFiveRFPermutationCILow}{0.034}
\renewcommand{\LBMFiveRFPermutationCIHigh}{0.054}
\renewcommand{\LBMFiveDiagnosticMin}{0.036}
\renewcommand{\LBMFiveDiagnosticMax}{0.084}
}{}

\title{LeakBench-Tab: Separating Contamination Validity, Statistical Detectability, and Model Exploitability in Tabular Learning}
\author{Anonymous Submission}
\affiliations{}

\begin{document}

\maketitle

\ifLBResultsReady\else
\begin{center}
\fbox{\parbox{0.94\linewidth}{\centering\bfseries
INTERNAL DRAFT --- RESULTS BLOCKED.\\
The complete confirmatory \texttt{paper\_claims.json} has not passed the
fail-closed paper-number generator.  No empirical value in this PDF is
submission-ready.}}
\end{center}
\fi

\begin{abstract}
Prediction-time feature contamination can invalidate tabular model evaluation,
yet an invalid feature need not be statistically conspicuous or useful to a
particular learner.  We present LeakBench-Tab, a measurement benchmark that
separates prediction-time validity (C), statistical detectability (D), and
model exploitability (X).  Its frozen controlled design crosses
\LBControlledTaskCount{} panel tasks, \LBMechanismCount{} mechanisms,
\LBStrengthCount{} strengths, \LBCoreModelCount{} model families, and
\LBSeedCount{} injection seeds, for \LBExpectedCells{} paired model-training
cells.  Mechanism-level construction invariants enforce the declared prediction
boundary, and hierarchical inference treats the task rather than the cell as
the independent unit.  A separate \LBExpectedDiagnosticCells{}-cell sensitivity
matrix compares \LBDiagnosticMethodCount{} oracle-blind, pre-test development
rankers; mutual information is the frozen primary diagnostic.
\ifLBResultsReady
The final release contains \LBSuccessfulCells{} successful cells
(\LBCompletionRate\% completion).  The declared simple-minus-structured
harm contrast is \LBSimpleStructuredStatus{} ($\Delta=
\LBSimpleStructuredDifference$, 95\% CI
$[\LBSimpleStructuredCILow,\LBSimpleStructuredCIHigh]$), while M03, M08, and
M09 expose mechanism-level profiles that cannot be summarized by one leakage
severity score.  Five fixed real-data case studies provide bounded consistency
checks only; they are not used for population-level or
zero-shot-transfer claims.
\else
Empirical findings are intentionally withheld until the complete confirmatory
release passes the frozen integrity and claim-status checks.
\fi
\end{abstract}

\section{Introduction}
Prediction-time leakage arises when a model is trained or evaluated with
information that would not be valid at its declared decision boundary.  Direct
copies of the target are obvious, but practical contamination also enters
through post-outcome fields, future-window aggregates, and entity or source
statistics whose legitimacy depends on lineage and timing
\citep{kaufman2011leakage,kapoor2023leakage,davis2024framework}.  Such fields can
make offline performance optimistic and comparisons between models misleading.

Association alone does not establish invalidity.  Conversely, an invalid field
may be hard to locate from samples alone, and a near-zero performance change
does not make that field legitimate.  We therefore separate three questions:
\emph{contamination validity} (C), determined by the prediction boundary;
\emph{statistical detectability} (D), measured as localization quality; and
\emph{model exploitability} (X), measured by paired permissive-versus-strict
performance inflation.

We ask: \textbf{RQ1:} Which prediction-time contamination mechanisms can
oracle-blind, pre-test development diagnostics locate, and how
diagnostic-conditional is that conclusion?  \textbf{RQ2:} Which mechanisms are
exploited across linear, tree-ensemble, boosting, and neural tabular learners?
\textbf{RQ3:} Which C/D/X profiles persist across mechanism strengths, model
families, and explicitly bounded real-data case studies?

Our contribution is a benchmark-level operationalization rather than a new
definition of leakage or a new detector.  First, LeakBench-Tab records C, D, and
X as distinct estimands tied to a declared prediction boundary.  Second, it
provides \LBMechanismCount{} construction-audited mechanisms on
\LBControlledTaskCount{} deterministic controlled panel tasks, with paired
strict and permissive evaluation.  Third, it freezes exact task-level inference
for the directional category contrast and task-hierarchical uncertainty
analyses for descriptive mechanism, diagnostic, and model-family summaries.
Fourth, it supplements
controlled evidence with
\LBNaturalTaskCount{} fixed real-data case studies whose scope is explicitly
limited to bounded consistency checks.  Metadata localization and governance
interventions are outside the confirmatory claim set reported here.

\section{Related Work}
\label{sec:related}
\paragraph{Leakage semantics and scientific validity.}
Prediction-time legitimacy and the ``no time machine'' principle precede this
work \citep{kaufman2011leakage}.  Later analyses connect leakage to failures of
scientific reproducibility and domain-specific label leakage
\citep{kapoor2023leakage,davis2024framework}.  Controlled neuroimaging audits
also show that different leakage forms can produce large or minor performance
inflation \citep{rosenblatt2024connectome}.  Most closely, a recent tabular
study applies 29 workflow perturbations across 2,047 i.i.d. and 129 temporal
datasets to compare leakage-type effect sizes \citep{roth2026leakage}.  That
work provides substantially broader dataset coverage; our complementary
resolution is within-task feature localization, construction-audited
prediction boundaries, and paired diagnostic/model profiles for eleven fixed
mechanisms.
LeakBench-Tab does not redefine
these concepts or claim that variable inflation is new.  Its distinct scope is
a domain-agnostic tabular mechanism registry that jointly evaluates semantic
validity, feature-level localization, and paired inflation across diagnostic
and downstream model families.

\paragraph{Tabular model evaluation.}
Broad comparisons such as TableShift and TabZilla study distribution shift and
model-family performance across tabular tasks
\citep{gardner2023tableshift,mcelfresh2023tabzilla}.  Other work establishes
strong tree baselines and modern neural alternatives
\citep{grinsztajn2022trees,gorishniy2025tabm}.  Our axis of comparison is
different: every learner receives paired views of the same frozen injected
task, so the estimand is contamination harm rather than unpaired predictive
rank.

\paragraph{Diagnostics versus exploitation.}
The distinction between contamination, memorization, and downstream
exploitation is not new: \citet{magar2022contamination} showed in NLP
pretraining that memorized contaminated examples need not be exploited after
fine-tuning.  Our use of ``exploitability'' therefore claims no conceptual
priority.  We operationalize a different setting: prediction-time tabular
feature contamination C, localization of encoded contaminated columns D, and
paired strict-versus-permissive degradation X.  Static program analysis offers
a complementary route by detecting leakage-prone data-science code before
execution \citep{subotic2022static}; our D diagnostics instead localize
contaminated columns in a frozen task.
Feature-level precision--recall summaries are appropriate when contaminated
fields are sparse \citep{saito2015precision}, but successful localization does
not imply that every learner will exploit the field.  We use mutual information
as the primary D measure and compare it with absolute correlation, logistic
coefficient magnitude, and random-forest permutation importance.  All four are
oracle-blind and pre-test: the first three use training rows only, while the
forest is fitted on training rows and permutation importance is evaluated with
the frozen validation labels.  Their disagreement is descriptive sensitivity,
not a procedure for selecting the best ranker after seeing contamination
labels.  X is reported separately for each downstream model.  The benchmark is measurement
infrastructure, not a proposal for a universally superior leakage detector or
governance policy.

\section{Problem Setup and the C/D/X Framework}
\label{sec:setup}
Let $\mathcal{D}=\{(\mathbf{x}_i,y_i)\}_{i=1}^{n}$ be a binary tabular task
with declared prediction time $t_p$.  A field may be invalid because it becomes
available after $t_p$, encodes future or outcome information, or violates a
declared lineage or population boundary.  The \emph{strict} view removes all
mechanism-labeled contaminated fields; the paired \emph{permissive} view retains
them.  All train, validation, and test row assignments remain fixed between the
two views.

\subsection{Contamination Validity (C)}
For field $j$, define
\begin{equation}
C_j = \mathbf{1}\!\left[j\text{ is invalid under the declared boundary at }t_p\right].
\end{equation}
$C_j$ is fixed by task semantics and the mechanism contract.  It is not inferred
from mutual information, model importance, attribution, or downstream
performance.

\subsection{Statistical Detectability (D)}
An oracle-blind diagnostic assigns a score $s_j$ to every encoded column
without access to the contamination mask or test rows.
Localization is evaluated against $C_j$ with feature-level average precision
(AP), normalized against contaminated-field prevalence $\pi$ as
\begin{equation}
D = \frac{\mathrm{AP}-\pi}{1-\pi}.
\end{equation}
We also retain top-$5$ recall.  Mutual information is the frozen primary
ranker used in the headline C/D/X profiles and uses the same fixed
\texttt{random\_state=42} for every injected task.  A sensitivity suite additionally
uses absolute correlation, logistic coefficient magnitude, and random-forest
permutation importance.  MI, correlation, and logistic coefficients use
training rows only; the forest is trained on training rows and scored by
permutation on the frozen validation split.  Each method's D is computed once
per injected task and is therefore invariant to the downstream exploitability
model.  D is diagnostic- and representation-conditional; no per-task
best-method maximum is a deployable selector.

\subsection{Model Exploitability (X)}
For learner $m$ and a metric for which larger is better, paired harm is
\begin{equation}
X_m = q_m(\mathcal{D}_{\mathrm{permissive}})
      -q_m(\mathcal{D}_{\mathrm{strict}}).
\end{equation}
We use test AUROC as $q$.  Positive X denotes measured performance inflation;
near-zero X denotes no material benefit under this protocol.  Negative X is not
evidence that a learner recognizes or rejects invalidity.

\begin{figure}[t]
\centering
\small
\setlength{\fboxsep}{4pt}
\fbox{\parbox{0.91\columnwidth}{\centering
\textbf{Declared prediction-time information boundary}\\
Which fields and lineage are available when the prediction is made?}}
\par\vspace{-0.15em}$\downarrow$\par\vspace{-0.15em}
\fbox{\parbox{0.91\columnwidth}{\centering
\textbf{C --- Contamination validity}\\
Semantic construction label; not inferred from association or performance.}}
\par\vspace{-0.15em}$\swarrow\hspace{0.36\columnwidth}\searrow$\par\vspace{-0.15em}
\begin{tabular}{@{}cc@{}}
\fbox{\parbox[t]{0.41\columnwidth}{\centering
\textbf{D --- Detectability}\\
Can an oracle-blind diagnostic localize the encoded contaminated columns?}} &
\fbox{\parbox[t]{0.41\columnwidth}{\centering
\textbf{X --- Exploitability}\\
How much does the paired permissive view inflate held-out AUROC?}}
\end{tabular}
\caption{C is fixed by the declared information boundary.  D and X are
separate empirical measurements conditional on a diagnostic, representation,
and learner; neither determines C or the other empirical axis.}
\label{fig:cdx-framework}
\end{figure}

\section{LeakBench-Tab Design}
\label{sec:design}
\subsection{Controlled Tasks and Mechanisms}
The controlled registry contains \LBControlledTaskCount{} deterministic panel
tasks spanning five data-generating archetypes.  Each task includes ordered
timestamps and sufficient entity and source support for the structured
mechanisms.  Generator seeds are separated from injection seeds, and every
mechanism-strength-seed tuple has a content hash and a fixed chronological
60/20/20 split.

Table~\ref{tab:mechanisms} groups mechanisms by construction rather than by
observed performance.  Strict-future mechanisms exclude the current row; M08
uses future entity outcomes with shrinkage and a leave-one-out prior; M09 uses
outcome-dependent source collection with a complete one-hot representation and
a random permutation of source category IDs to prevent ordinal encoding; and
M10 keeps its legitimate component exactly equal
to a clean input while labeling the contaminated component separately.  These
invariants make C auditable before any model is fitted.
M09's D localizes eight encoded one-hot columns representing one source
variable, so its detectability is representation-conditional rather than a
field-level generality claim across encodings.

\begin{table*}[t]
\centering
\small
\begin{tabular}{llll}
\toprule
ID & Mechanism & Category & Invalidity source / structure \\
\midrule
M01 & Direct target copy & Simple & Target-derived global field \\
M02 & Noisy target proxy & Simple & Noisy target-derived global field \\
M06 & Redundant leakage cluster & Simple & Correlated target proxies \\
M10 & Mixed leakage & Simple & Legitimate and contaminated components \\
M03 & Nonlinear target interaction & Boundary & Nonlinear target-derived interaction \\
M07 & Sparse subgroup leakage & Boundary & Contamination restricted to a subgroup \\
M11 & Graph-mediated leakage & Boundary & Connected contamination component \\
M04 & Post-outcome aggregation & Structured & Strictly post-outcome temporal aggregate \\
M05 & Temporal look-ahead & Structured & Strictly future-window information \\
M08 & Entity-conditioned leakage & Structured & Future entity-outcome statistic \\
M09 & Source-conditioned leakage & Structured & Outcome-dependent source assignment \\
\bottomrule
\end{tabular}
\caption{Frozen mechanism taxonomy.  Categories are declared by construction
and are not reassigned after observing model performance.}
\label{tab:mechanisms}
\end{table*}

\subsection{Models}
The five core learners are logistic regression, random forest
\citep{breiman2001random}, CatBoost \citep{prokhorenkova2018catboost}, LightGBM
\citep{ke2017lightgbm}, and the official TabM implementation
\citep{gorishniy2025tabm}.  Scaling, preprocessing, and diagnostics are fit on
training data only.  The TabM tasks are serialized locally and verified by task,
split, and bundle hashes before GPU fitting; no surrogate implementation or
silent model fallback is accepted.

\subsection{Real-Data Case Studies}
Five fixed case studies---UCI Bank Marketing \citep{uciBankMarketing}, a local
Lending Club accepted-loan mirror, U.S. flight delays \citep{btsOnTime}, Chicago
food inspections \citep{chicagoFoodInspections}, and NYC 311
\citep{nyc311}---declare a prediction boundary and audit-confirmed post-boundary
fields before evaluation.  They use the four CPU core models.  These studies
fit categorical vocabularies on training rows only, map unseen validation/test
levels to a reserved unknown value, and exclude globally encoded date strings.
The earlier full-table-vocabulary run is superseded and cannot source claims.
These studies
ask whether paired harm appears under specific real task boundaries; they do
not train a cross-task detector, test zero-shot metadata transfer, or justify
population inference from five selected tasks.  The Lending mirror is treated
as local-only and non-redistributable until its source and license are verified.

\section{Experimental Protocol}
\label{sec:protocol}
The frozen controlled matrix crosses \LBControlledTaskCount{} tasks,
\LBMechanismCount{} mechanisms, \LBStrengthCount{} strengths,
\LBCoreModelCount{} learners, and \LBSeedCount{} injection seeds, yielding
\LBExpectedCells{} paired model-training cells.  The primary estimand is mean
paired AUROC harm.  Clean-model results are cached only within an identical
task/model/seed key, and every failed cell is retained rather than silently
dropped.

The task is the independent unit.  Confirmatory intervals use a 20,000-draw
hierarchical bootstrap that resamples tasks and then injection seeds.  Mechanism
tests use task-level sign flips with Holm correction over the declared family;
the three category contrasts use exact task-level sign flips and form a separate
Holm family.  M08 entity sensitivity shares each entity draw across all models
and strengths within a task/seed group before outer resampling.  M09 shares each
source-category draw across models within a task/seed/strength group, but its
eight source IDs are constructed categories rather than sampled population
clusters; that interval is descriptive designed-category reweighting only.
Because the 20 tasks are a designed registry rather than a random
sample from a defined task population, these intervals quantify stability to
registry reweighting; they do not support population inference to arbitrary
tabular tasks.  Association between mean D and X is explicitly descriptive: we
report global rank correlation, category-only and category-plus-D fits,
within-category permutation, and leave-one-mechanism-out (LOMO) diagnostics.
The separate sensitivity matrix contains
\LBControlledTaskCount{}$\times$\LBMechanismCount{}$\times$
\LBStrengthCount{}$\times$\LBSeedCount{}$\times$
\LBDiagnosticMethodCount{}$=\LBExpectedDiagnosticCells{}$ diagnostic cells on
the same immutable tasks.  It is complete only when all
\LBSuccessfulDiagnosticCells{} cells pass integrity checks; its rows are not
added to the \LBExpectedCells{} model-training count.

Post-unblinding evidence reconciliation found that the original sensitivity
runner had used the injection seed for its mutual-information rows, while the
headline analysis and frozen task manifest used the fixed seed above.  A
scope-locked amendment therefore replaces all and only the 5,500 MI rows with
the already-frozen task-manifest localization metrics; the other 16,500 rows
and the method set are unchanged.  The raw suite remains in the
archive as provenance but is not paper-facing evidence.

A second post-unblinding methodological audit found that bootstrap tail areas
had been labeled as category-contrast $p$-values, D--X intervals had resampled X
while holding D fixed, and cluster draws had been made independently per model
cell.  A second audit then found that the initial cluster amendment still drew
M08 entities independently across five injection seeds that reuse the same test
rows, labels, and entities.  The final scope-locked amendment replaces those summaries:
exact task sign flips with Holm adjustment are used for category tests, D and X
are jointly resampled, and cluster/category draws are synchronized at their
shared task level.  For M08, one entity draw is shared across all five seeds,
models, and strengths within a dataset while seed-specific effects are retained
for outer resampling.  Prediction arrays and task hashes are checked directly
against the frozen bundles.  Mechanism-profile comparisons remain descriptive
because no decision thresholds for them were bound by the original protocol;
both superseded cluster summaries are barred from claim provenance.

\ifLBResultsReady
The canonical manifest contains \LBSuccessfulCells{}/\LBExpectedCells{}
successful cells (\LBCompletionRate\%).
\else
Completion is not asserted in this draft because the final confirmatory claim
release is absent.
\fi

\section{Results}
\label{sec:results}
\ifLBResultsReady
\subsection{The Frozen Category Contrast}
The simple-minus-structured mean harm contrast was
$\LBSimpleStructuredDifference$ (95\% hierarchical bootstrap CI
$[\LBSimpleStructuredCILow,\LBSimpleStructuredCIHigh]$; Holm-adjusted
$p=\LBSimpleStructuredHolmP$ from an exact two-sided task-level sign-flip test).
\ifLBSimpleStructuredSupported
This satisfies the declared release criterion.
\else
This does not satisfy the declared release criterion and is reported as a
negative confirmatory result.
\fi
The direction was positive for \LBModelPositiveDirectionCount{} of five
learners; \LBModelCIExcludesZeroCount{} model-specific intervals excluded zero.
These model-specific counts are descriptive and do not enter the support rule.
This category average is not interpreted as a universal separation because the
structured set contains a declared counterexample.

\subsection{Crossed Mechanism Profiles}
For M03, normalized AP was \LBMThreeDetectability{} (95\% CI
$[\LBMThreeDetectabilityCILow,\LBMThreeDetectabilityCIHigh]$) and paired harm
was \LBMThreeHarm{} (95\% CI
$[\LBMThreeHarmCILow,\LBMThreeHarmCIHigh]$).  This profile is descriptive;
we do not apply a post hoc low-D or high-X threshold.
For M08, normalized AP was \LBMEightDetectability{} (95\% CI
$[\LBMEightDetectabilityCILow,\LBMEightDetectabilityCIHigh]$) and paired harm
was \LBMEightHarm{} (95\% CI
$[\LBMEightHarmCILow,\LBMEightHarmCIHigh]$); the synchronized entity-cluster sensitivity CI
was $[\LBMEightClusterCILow,\LBMEightClusterCIHigh]$.
This interval is descriptive uncertainty under registry reweighting and is not
an equivalence or practical-null test.

Within this fixed registry and encoded source representation, M09 is a
descriptive structured counterexample to a uniform structured-harmlessness
pattern.  Its normalized AP was
\LBMNineDetectability{} (95\% CI
$[\LBMNineDetectabilityCILow,\LBMNineDetectabilityCIHigh]$) and paired harm was
\LBMNineHarm{} (95\% CI
$[\LBMNineHarmCILow,\LBMNineHarmCIHigh]$).
Here D is encoded-column localization over an eight-column one-hot
representation of one semantic source field; it must not be generalized to
field-level detectability under other representations.  Reweighting the eight
constructed source categories gives the descriptive interval
$[\LBMNineReweightingLow,\LBMNineReweightingHigh]$; it is not a source-population
cluster confidence interval and is not a support criterion for this claim.
It does not establish a population-level source effect or transfer to another
source encoding.

\begin{figure}[t]
\centering
\IfFileExists{\LBCDXScatterPath}{%
\includegraphics[width=\columnwidth]{\LBCDXScatterPath}}{%
\PackageError{LeakBench-Tab}{Missing final C-D-X scatter}{Generate it from the
canonical confirmatory summaries; pilot figures are forbidden.}}
\caption{Mechanism-level primary-MI detectability (D) and paired harm (X), with
task-hierarchical intervals and construction categories.  M03, M08, and M09 are
descriptive registry profiles; no post hoc profile threshold determines their
status.}
\label{fig:cdx-scatter}
\end{figure}

\subsection{Model-Family Heterogeneity}
The paired design also exposes learner-conditional exploitation: the same
construction label and diagnostic score can accompany different harm across
linear, bagged-tree, boosted-tree, and neural learners.  We report these
differences as model-family summaries, not causal capacity effects.

\begin{figure}[t]
\centering
\IfFileExists{\LBMechanismModelHeatmapPath}{%
\includegraphics[width=\columnwidth]{\LBMechanismModelHeatmapPath}}{%
\PackageError{LeakBench-Tab}{Missing final mechanism-model heatmap}{Generate it
from the canonical frozen table; pilot figures are forbidden.}}
\caption{Paired harm by mechanism and model family under a common scale.  All
55 fixed mechanism--model cells are retained; the contrasts are descriptive.}
\label{fig:mechanism-model}
\end{figure}

\subsection{Detectability Is Diagnostic-Conditional}
Across the four oracle-blind, pre-test development rankers, normalized AP ranged from
\LBMThreeDiagnosticMin{} to \LBMThreeDiagnosticMax{} for M03, from
\LBMFourDiagnosticMin{} to \LBMFourDiagnosticMax{} for M04, and from
\LBMFiveDiagnosticMin{} to \LBMFiveDiagnosticMax{} for M05.  These ranges are
reported as descriptive method sensitivity around the primary MI analysis.
Without paired simultaneous method-comparison intervals, we do not attach an
inferential detector-ranking claim.  We do not
select a ranker per mechanism and do not interpret the maximum as an operational
detector.

\begin{figure}[t]
\centering
\IfFileExists{\LBStrengthDiagnosticPath}{%
\includegraphics[width=\columnwidth]{\LBStrengthDiagnosticPath}}{%
\PackageError{LeakBench-Tab}{Missing final strength/diagnostic figure}{Generate
it from the complete confirmatory diagnostic matrix; pilot figures are
forbidden.}}
\caption{Strength-response and descriptive diagnostic sensitivity under four
frozen oracle-blind, pre-test development rankers.  The first three use
training rows only; RF permutation is evaluated on the frozen validation
split.  Mutual information remains primary; the best observed
ranker per mechanism is not a deployment rule.}
\label{fig:diagnostic-robustness}
\end{figure}

\subsection{D--X Association Is Descriptive, Not a Transfer Rule}
Across the eleven mechanism means, Spearman $\rho$ was \LBGlobalSpearman{}
(95\% joint-paired D--X bootstrap interval
$[\LBGlobalSpearmanCILow,\LBGlobalSpearmanCIHigh]$).  The same task/seed draws
resample both axes.  Category alone
had in-sample $R^2=\LBCategoryRSquared$; adding D yielded
$R^2=\LBCategoryPlusDRSquared$ (increment \LBIncrementalRSquared{},
within-category permutation $p=\LBIncrementalPermutationP$).  The corresponding
LOMO values were \LBCategoryLomoRSquared{} and
\LBCategoryPlusDLomoRSquared{} (increment \LBIncrementalLomoRSquared{}).
These eleven points describe this mechanism registry; they do not establish a
zero-shot predictor of exploitability.

\subsection{Fixed Real-Data Boundary Checks}
Across \LBNaturalTaskCount{} fixed case studies and \LBNaturalModelCount{} CPU
learners, whether every task-level effect was positive was
\LBNaturalAllPositive{}.  Mean paired harm was \LBNaturalMeanHarm{} (descriptive
task-bootstrap interval $[\LBNaturalCILow,\LBNaturalCIHigh]$; exact task-level
sign-flip $p=\LBNaturalSignFlipP$).
These are descriptive results for the five fixed cases.  Neither the interval
nor the five-task sign test is used as a support decision or population claim.

\begin{table}[t]
\centering
\small
\begin{tabular}{@{}p{0.67\columnwidth}p{0.25\columnwidth}@{}}
\toprule
Paper-facing evidence & Claim scope \\
\midrule
Simple minus structured: $\Delta=\LBSimpleStructuredDifference$,
95\% CI $[\LBSimpleStructuredCILow,\LBSimpleStructuredCIHigh]$,
Holm $p=\LBSimpleStructuredHolmP$ & \LBSimpleStructuredStatus \\
M03, M08, and M09 crossed C/D/X profiles & descriptive only \\
D--X association over eleven fixed mechanisms & \LBRelationStatus \\
Five selected real-data boundary checks & \LBNaturalStatus \\
\bottomrule
\end{tabular}
\caption{Claim scope is part of the result: only the pre-existing directional
category contrast receives a binary support decision.}
\label{tab:claim-scope}
\end{table}
\else
\paragraph{Blocked draft.}
No confirmatory result is rendered because
\texttt{generated/result\_macros.tex} has not been produced from a complete,
confirmatory, non-pending \texttt{paper\_claims.json}.  Pilot estimates are
deliberately excluded from the manuscript.
\fi

\section{Limitations and Conclusion}
\label{sec:conclusion}
The controlled tasks are synthetic panel generators, not twenty real datasets.
The registry covers eleven designed mechanisms and four oracle-blind, pre-test
development rankers; neither set is exhaustive.  AUROC harm addresses binary discrimination
and does not replace calibration, decision-utility, or deployment monitoring.
The five real-data audits are selected, boundary-specific case studies with
sparse contamination labels, so their uncertainty intervals are descriptive
and their conclusions do not generalize to an unspecified task population.
Model coverage is limited to five core families, and differences between
families are associational rather than causal.  Metadata localization and
governance comparisons remain outside this paper's confirmatory evidence.

\ifLBResultsReady
Within these limits, LeakBench-Tab shows why validity, detectability, and
exploitability should be reported separately: validity is semantic, D and X are
empirical, and mechanism-level crossed profiles prevent either empirical axis
from serving as a universal leakage-severity score.
\else
The non-empirical conclusion is limited to the benchmark design; the final
empirical conclusion remains blocked pending the confirmatory release.
\fi

\FloatBarrier
\label{lb:last-main-content-page}
\bibliography{references}

\end{document}
